% =====================================================================
%  sections/sec3.tex —— 3 Hitchin 方程的非正则奇点（原文 pp.5–12）
% =====================================================================
\section{Hitchin 方程的非正则奇点}
\label{sec:hitchin}

AD 理论由六维 $A_1$ 理论在带非正则奇点的黎曼面上紧致化构造出来 \cite{3}。这个描述是这样得到的：先对已知 UV 理论的 Seiberg--Witten 曲线取标度极限，再把所得的 Seiberg--Witten 曲线映成 Hitchin 系统的描述。把取标度极限的做法推广到更高秩的理论相当困难，而且标度极限本身还有若干微妙之处，近期文献 \cite{22} 讨论过。我们的做法是：直接从 Hitchin 方程的非正则奇异解出发，用非正则奇点的分类来分类四维理论。同样的思路在正则穿刺点的情形似乎也运行良好——那里 Hitchin 方程的解确实有杨图分类，与从物理出发的推导一致 \cite{23}。

\routenote{原文 p.5。这里作者明确换了路线：不再走「对已知 UV 理论取标度极限」这条路，而是把 Hitchin 方程的非正则奇点当作出发点，用奇点的分类来分类理论。后面第 3.1 节的全部分类工作都是在为这条路线铺路。}

AD 理论的新特征就是分数标度维与有量纲的耦合常数，而只用正则奇点无法容纳这两个新特征。引入非正则性恰好解决了问题：第一，黎曼面上的坐标 $z$ 在 $U(1)_R$ 对称性下按非平凡方式变换，于是有了分数标度；第二，高阶极点中的系数就是有量纲的耦合常数。几何描述与物理量之间还有许多别的奇妙对应，后面会详细讨论。本节主要从 Hitchin 方程的角度描述非正则奇点，想直接看物理应用的读者初读时可以跳过本节。

下面先把物理量与 Hitchin 方程的几何侧面之间的各种对应关系细看一遍，其正则穿刺点的情形稍后要推广到非正则奇点的情形。文献 \cite{2} 构造的四维 $\sN{=}2$ 理论，是把六维 $(2,0)$ SCFT 在带正则穿刺点（含拓扑扭曲）的黎曼面上紧致化得到的。Hitchin 方程同样定义在这个黎曼面上，而穿刺点对应于 Hitchin 方程的正则奇异解。边界条件固定的 Hitchin 方程解模空间是超凯勒流形，它被认同为四维理论在 $\R^3\times S^1$ 上的库仑分支，其中 $S^1$ 是一个圆。事实上，Hitchin 模空间就是五维极大超杨–米尔斯理论在带穿刺黎曼面上紧致化后的希格斯分支，而它是原三维理论的镜像 \cite{24}。由于非重整定理，希格斯分支没有量子修正；正因如此，Hitchin 方程的经典图像才能编码原理论库仑分支的信息。

几何参数是带穿刺黎曼面的复结构模，以及一阶极点的系数；它们分别被认同为规范耦合常数与质量参数。场论的红外行为编码在模空间里，而后者是超凯勒流形，其复结构由 $\CP^1$ 参数化。在其中一个复结构 $I$ 下，Hitchin 模空间是一个可积系统，其谱曲线被认同为 Seiberg--Witten 曲线。在复结构 $J$ 下，模空间上的每一点参数化黎曼面上的一个平坦联络，而这个描述对非正则奇点的分类有用。取定黎曼面上的一个复结构后，Hitchin 方程写成
\begin{equation}
F - \phi\wedge\phi = 0,\qquad D\phi = D\ast\phi = 0,
\label{eq:3.1}\tag{3.1}
\end{equation}
其中 $A$ 是联络，$\phi$ 是一个一形式，称为 Higgs 场 \cite{8,25}。若记 $\sA = A + \ii\phi$，则 Hitchin 方程蕴含 $\sA$ 的曲率平坦。事实上还可以引入谱参数，定义一族依赖于谱参数的平坦联络。平坦联络 $\sA$ 绕奇点的单值性，可通过解下面的平坦截面方程算出：
\begin{equation}
(\partial_z + \sA_z)\,\psi = 0,
\label{eq:3.2}\tag{3.2}
\end{equation}
这在局部不过就是圆盘上的一个一阶微分方程。

规范群为 $\SU(N)$ 时，Hitchin 方程最简单的非正则奇异解已在文献 \cite{26} 中研究过，奇点附近的场具有下面的形式
\begin{equation}
\phi = \frac{u_n}{z^n} + \cdots + \frac{u_2}{z^2} + \frac{u_1}{z} + \text{c.c} + \cdots,\qquad
A = \alpha\,\dd\theta .
\label{eq:3.3}\tag{3.3}
\end{equation}
这里取局部坐标 $z = re^{\ii\theta}$，$u_1,\dots,u_n$ 都是特征值互异的对角矩阵（这可以借助规范对称性做到）。省略号代表正则项，下面的公式里会略去，但它们始终存在。这种把平坦联络 abelianization 的做法，对我们寻找 Hitchin 方程的非正则奇异解至关重要。

\zhnote{[来源：原文 p.6 与 p.5 的第一条脚注；I：本库无先例。] abelianization 指把取值在 $\mathfrak{sl}(N)$ 中的联络通过规范变换对角化到 Cartan 子代数里，即把非阿贝尔对象拆成一族阿贝尔的对象。中文学界另有一个「阿贝尔化」，专指群论里的 $G\to G/[G,G]$；两者所指不同，混用会误导，故保留英文。}

存在非正则奇点时，Hitchin 模空间仍是超凯勒流形 \cite{10}，其性质与正则奇点的情形有许多共同之处。特别是：复结构 $I$ 线性依赖于一阶极点的系数，而后者应被认同为质量参数 \cite{26}；复结构 $I$ 定义一个可积系统，其谱曲线为
\begin{equation}
\det\bigl(x - \Phi(z)\bigr) = 0,
\label{eq:3.4}\tag{3.4}
\end{equation}
其中 $\Phi$ 是 $\phi$ 的全纯部分，而这个谱曲线被认同为 Seiberg--Witten 曲线。使每一点代表一个平坦联络的那个复结构 $J$ 对分类有用，下面我们专注于它。对 (3.3) 给出的非正则奇点，平坦联络的全纯部分具有下面的形式
\begin{equation}
\sA_z = \frac{u_n}{z^n} + \frac{u_{n-1}}{z^{n-1}} + \cdots \frac{u_2}{z^2} + \frac{u'_1}{z}.
\label{eq:3.5}\tag{3.5}
\end{equation}

微分方程 (3.2) 有一种有趣的 Stokes 现象，它对定义绕非正则奇点的单值性很重要。为了完备起见，我们把 Stokes 现象的若干侧面回顾一遍；感兴趣的读者可在 \cite{26,27} 中找到更多细节。这些 Stokes 矩阵的出现，源自方程 $(\partial + \sA_z)\psi = 0$ 解的渐近行为。先设规范群为 $U(1)$，微分方程就成为
\begin{equation}
\frac{\dd\psi}{\dd z} = -\Bigl(\frac{q_n}{z^n} + \frac{q_{n-1}}{z^{n-1}} \cdots + \frac{q_1}{z} + B(z)\Bigr)\psi,
\label{eq:3.6}\tag{3.6}
\end{equation}
这里 $B(z)$ 是 $z=0$ 处正则的全纯函数。解非常简单：
\begin{equation}
\psi = c(z)\exp Q(z),\qquad
Q(z) = \frac{q_n}{(n-1)z^{n-1}} + \frac{q_{n-1}}{(n-2)z^{n-2}} \cdots + q_1(-\ln z),
\label{eq:3.7}\tag{3.7}
\end{equation}
这里 $c(z)$ 是在奇点附近不收敛的形式幂级数。这个解是更高秩情形的解的构件：后者不过是上述解构成的向量，带指标 $i = 1,2,\dots,n$。然而，沿通往奇点的不同路径，解向量的各分量有不同的渐近行为，因为
\begin{equation}
\Bigl|\frac{\exp Q^i(z)}{\exp Q^j(z)}\Bigr| \to
\Bigl|\exp\Bigl(\frac{q_n^i - q_n^j}{(n-1)z^{n-1}}\Bigr)\Bigr|,\qquad z\to 0 .
\label{eq:3.8}\tag{3.8}
\end{equation}
渐近行为取决于 $\operatorname{Re}\bigl((q_n^i-q_n^j)/((n-1)z^{n-1})\bigr)$ 的符号。在不同角区，符号不同；容易看出共有 $2(n-1)$ 个角区，它们被称为 Stokes 扇区，Stokes 扇区的边界被称为 Stokes 射线。

要点在于：若某个区域内没有 Stokes 射线，则在该区域内具有给定渐近行为的解并不唯一。例如，若在该区域内 $|\exp(Q_i(z))/\exp(Q_j(z))|\gg 0$，那么解 $\psi'_i(z) = \psi_i(z) + \lambda\psi_j(z)$ 与 $\psi_i$ 有相同的渐近行为。若该角区内有 Stokes 射线，这种自由度就不存在了：被 Stokes 射线分开的两个区域渐近行为不同，因此只要该区域内有 Stokes 射线，区域内 $\psi_i(z)$ 与 $\psi_j(z)$ 这两组解就被固定住了。

取一个角宽为 $\pi/(n-1)$、边界不是 Stokes 射线的区域。把这个区域按 $\pi/(n-1)$ 的整数倍旋转，就得到圆盘的一个覆盖。可以把每个扇区稍稍放大，使得重叠区域内没有 Stokes 射线。在每个扇区内，对解向量的任意一对分量都存在一条 Stokes 射线，而整个解由该区域内的渐近行为唯一确定。然而在重叠区域上没有 Stokes 射线，两个区域中的两套解由一个对角元全为一的上三角矩阵相联系。这个矩阵就叫 Stokes 矩阵，它随后被用来构造单值性矩阵。

总共有 $2(n-1)$ 个 Stokes 矩阵，它们的乘积定义了广义单值性的一部分。单值性还有一个来自正则奇异项的贡献，该项贡献 $N-1$ 个参数。最后，我们还需减去 $T_c$ 群（其维数为 $N-1$）的贡献。于是，指定局部单值性所需参数的总数为
\begin{equation}
c_n = (n-1)(N^2 - N).
\label{eq:3.9}\tag{3.9}
\end{equation}
对一个绕过非正则奇点、取定基点的路径，其单值性还有另一份贡献，即额外 $(N^2-1)$ 个参数。然而，模空间是由指定领头阶系数来定义的，这给出 $-(N-1)$ 的贡献。于是非正则奇点对模空间的贡献总量为
\begin{equation}
d = n(N^2 - N).
\label{eq:3.10}\tag{3.10}
\end{equation}
若有不止一个非正则奇点，模空间的总维数是各局部贡献之和减去规范群的贡献：
\begin{equation}
\dim(\sM) = \sum_i d_i - 2(N^2 - 1).
\label{eq:3.11}\tag{3.11}
\end{equation}

\routenote{原文 pp.7--8。这两步计数（(3.9) 与 (3.10)）是本节最容易被读过去的地方，其逻辑是：先数「指定局部单值性需要多少参数」，再扣掉「由领头阶系数固定住的那部分」。注意 (3.10) 与 (3.9) 之差恰为 $(N^2-1)-(N-1)$，即绕奇点的路径给出的 $(N^2-1)$ 个参数减去领头阶系数固定掉的 $(N-1)$ 个。原文只写作 $T_c$ 群，未加解释；译文照录记号。}

非正则奇点对模空间的局部贡献，还可以用下述简单方式理解 \cite{28}：每个固定的正则半单矩阵定义李代数中的一个共轭类 $\sO_i$；微分算子 $d_{\sA}$ 把 $n$ 个共轭类之积中的一点，映成平坦联络模空间中的一点：
\begin{equation}
d_{\sA}: \sO_1\times\sO_2\times\sO_3\cdots\times\sO_n \to \sM .
\label{eq:3.12}\tag{3.12}
\end{equation}
由于 $\Sl(N)$ 的每个正则半单共轭类维数为 $N^2-N$，局部贡献的总维数就是 $d = n(N^2-N)$，与从 Stokes 矩阵得到的结论一致。

\subsection{非正则奇点的分类}
\label{sec:3.1}

本小节给出非正则奇点的分类。分类是借助指定高阶极点系数矩阵的形式来完成的。引入 $\omega = \tfrac{1}{z}$ 之后，方程 (3.2)\footnote{我们假定规范场系数与 Higgs 场的一阶系数属于同一个共轭类，由此把规范场对复平坦联络的贡献略去不计。} 变成
\begin{equation}
\frac{\dd\psi}{\dd\omega} = \Phi'(\omega)\psi,\qquad
\Phi'(\omega) = z^{-2}\Phi(1/z).
\label{eq:3.13}\tag{3.13}
\end{equation}
此时奇点位于 $\omega = \infty$。本文把奇点放在 $0$ 与放在 $\infty$ 交替使用：奇点放在 $\infty$ 便于分类、便于求 Seiberg--Witten 曲线，而放在 $0$ 便于考虑非正则奇点的碰撞。

关键在于：指定奇点行为的那组矩阵，可以借助形式规范变换化成对角形 \cite{26}；而这些对角矩阵正是我们要分类的对象。领头阶矩阵具有下面的普遍形式
\begin{equation}
\Phi(z) = \diag\bigl(z^{r_1-2}B_1,\ z^{r_2-2}B_2,\ \dots,\ z^{r_s-2}B_s\bigr),
\label{eq:3.14}\tag{3.14}
\end{equation}
其中 $B_i$ 是 $k_i\times k_i$ 的对角矩阵，其特征值的退化稍后再讨论；$r_1>r_2>\dots>r_s$ 是一组有理数，表示各 block 的极点阶。这里的关键之处是：所得的谱曲线在 $z$ 上只能出现整数幂，因此这些有理数 $r_i$ 必须具有下面的形式
\begin{equation}
r_i = n_i + \frac{j_i}{k_i},\qquad 0 < j_i \le k_i,
\label{eq:3.15}\tag{3.15}
\end{equation}
其中 $n_i$、$j_i$、$k_i$ 都是正整数。

下面研究各矩阵特征值的退化，只看 $B_1$ 这个 block 就够了（其他 block 类似）。若 $j_1\ne k_1$，特征值只能取下面的形式（这里假定 $j_1$ 与 $k_1$ 没有公因子；退化情形稍后再考虑）：
\begin{equation}
B_1 = \diag(1,\ \omega,\ \omega^2,\ \dots,\ \omega^{k_1-1}),\qquad j_1\ne k_1,
\label{eq:3.16}\tag{3.16}
\end{equation}
其中 $\omega = \exp\bigl(\tfrac{2\pi\ii}{k_1}\bigr)$（后文不再显写这个值，但只要 $k_1$ 选取得当，$\omega$ 总取上述形式）。这个特定形式来自「要求谱曲线在坐标 $z$ 上只有整数阶极点」。Higgs 场中出现分数幂的含义已在 \cite{26} 中解释过：奇点处会长出一条割线，需要做一次规范变换，把规范场从 $\Phi(\theta+2\pi)$ 变回 $\Phi(\theta)$。\footnote{也可以做规范变换，把上面的解变成具有整数极点的形式，但那对我们的目的不利，所以我们坚持上面这种异形表示。}这个规范变换本质上是特征值上的一个置换；因为绕奇点转一圈之后，新的解是
\begin{equation}
B'_1 = \diag\bigl(\omega^{j_1},\ \omega^{j_1+1},\ \omega^{j_1+2},\ \dots,\ \omega^{j_1+k_1-1}\bigr),\qquad j_1\ne k_1 .
\label{eq:3.17}\tag{3.17}
\end{equation}

由规范变换实现的置换，会把 Higgs 场带回原来的形式。次领头阶的矩阵则由「凡是与领头阶矩阵相容的形变都允许」这一要求确定：我们先前引入的那个置换，在绕奇点一圈之后必须把它变回原来的形式，因此下面这些项是允许的：
\begin{equation}
z^{m+\frac{l}{k_1}-2}\bigl(\omega^d,\ \omega^{d+1},\ \omega^{d+2},\ \dots,\ \omega^{k_1+d-1}\bigr),
\label{eq:3.18}\tag{3.18}
\end{equation}
其中 $m$、$l$ 是整数，$-1\le (m+\tfrac{l}{k_1}) < r_1$，而 $d$ 是只依赖于 $l$ 的一个固定整数。此外，下面这些项也是允许的：
\begin{equation}
z^{m-2}\,(a,\ a,\ \dots,\ a),
\label{eq:3.19}\tag{3.19}
\end{equation}
其中 $m<r_1$ 为整数——因为显然任何置换都会把它带回原形式。

若 $j_1 = k_1$，则极点阶 $r_i$ 是整数，$B_i$ 的特征值可以取任意退化形式
\begin{equation}
B_1 = \diag(a,\dots,a,\ b,\dots,b,\ c,\dots,c,\ \dots),\qquad j_1 = k_1,
\label{eq:3.20}\tag{3.20}
\end{equation}
这种退化可以用总格数为 $k_1$ 的一个杨图 $Y$ 记录下来，即第一列的高度等于取值为 $a$ 的特征值个数，其余依此类推。次领头阶的项同样只有整数阶极点，其特征值的退化由一个更大的杨图编码。例如，若领头阶的杨图是 $[3,2,1]$，意思是前三个特征值相同、接着两个特征值相同，等等；那么次领头阶的项可以用杨图 $[2,2,1,1]$ 表示，即前三个特征值现在拆成两块：其中两个仍然相同，第三个已经不同了。

领头阶矩阵有一种漂亮的图形表示，实践证明它对求 Seiberg--Witten 曲线非常有用：一个非正则奇点由二维整数格上的一个凸牛顿多边形表示。这个凸多边形的确定方式如下：先找一个点 $p_1$\footnote{$x$ 坐标要求小于 $N$，$z$ 坐标为正。}，使 $p_1$ 与点 $p_0 = (N,0)$ 的连线斜率为 $r_1-2$；再找另一个点 $p_2$\footnote{$p_i$ 的 $x$ 坐标小于 $p_{i-1}$，$z$ 坐标大于 $p_{i-1}$。}，使连线 $p_2p_1$ 的斜率为 $r_2-2$，依此类推。最终得到一个凸多边形，见图 \ref{fig:1}。若有某条线段的斜率是整数，上面的图形还不足以确定非正则奇点，此时还需要指定低阶极点矩阵的退化。下面是几个重要的例子：

\begin{enumerate}[label=\textup{\alph*.}]
  \item 上一小节所考虑的非正则奇点对应于 $r_1 = n$（$n$ 为整数），且只有一个 block；此外 $B_1$ 的特征值互不相同。
  \item 同样只有一个 block，但此时领头阶极点的阶 $r_1$ 可以是分数。
  \item 领头阶有两个 block，其形式为
  \begin{equation}
  \Phi = z^{r-2}\diag(0,\ 1,\ \omega,\ \dots,\ \omega^{N-2}).
  \label{eq:3.21}\tag{3.21}
  \end{equation}
  \item 领头阶只有一个 block，且 $r_1 = n$ 为整数。矩阵 $B_1$ 由杨图 $Y_n$ 指定，它表示 $B_1$ 特征值的退化。此时还需要一列杨图，使 $Y_n\subseteq Y_{n-1}\cdots\subseteq Y_1$，其中 $Y_{j-1}$ 是把 $Y_j$ 的每一列进一步分割得到的。
\end{enumerate}

注意这四类有一个共同点：领头阶矩阵本质上只有一个 block，而这些非正则奇点正是定义 AD 理论所需要的。此外，情形 a、b、c 的领头阶矩阵特征值互不相同。情形 a、b 可以用图 \ref{fig:2} 左图所示的牛顿多边形表示（情形 a 只需在边界上放适当的点），情形 c 则由右图表示。

\begin{figure}[tbp]
\centering
% Figure 1 -- Newton polygon of an irregular singularity (p_0, p_1, p_2)
\begin{tikzpicture}[x=1pt,y=1pt,scale=1.12,transform shape,>=stealth,line width=0.5pt]
  \def\u{14.145}
  % axes
  \draw[->] (0,0) -- (0,187.3);
  \draw[->] (0,0) -- (187.3,0);
  \node[inner sep=0pt] at (-2.2,202.5) {\fontsize{15}{18}\selectfont$z$};
  \node[inner sep=0pt] at (209.1,2.7)  {\fontsize{15}{18}\selectfont$x$};
  % lattice points on the boundary
  \foreach \i in {1,...,7} {\fill (0,{\i*\u}) circle[radius=2.7pt];}
  \foreach \i in {1,...,9} {\fill ({\i*\u},0) circle[radius=2.7pt];}
  % Newton polygon
  \draw (0,{7*\u}) -- ({3*\u},{6*\u}) -- ({6*\u},{4*\u}) -- ({9*\u},0);
  % labels of the boundary points
  \node[inner sep=0pt,font=\fontsize{10.6}{13}\selectfont] at (51.4,99.0)  {$p_2$};
  \node[inner sep=0pt,font=\fontsize{10.6}{13}\selectfont] at (106.3,62.7) {$p_1$};
  \node[inner sep=0pt,font=\fontsize{10.6}{13}\selectfont] at (138.1,16.8) {$p_0$};
\end{tikzpicture}

\caption{表示一个非正则奇点的牛顿多边形。边界上的每条线段代表一个 block，其斜率是该 block 的极点阶。}
\label{fig:1}
\end{figure}

\begin{figure}[tbp]
\centering
% Figure 2 -- left: one block (cases a,c); right: leading order partition [N-1,1] (case b)
\begin{tikzpicture}[x=1pt,y=1pt,>=stealth,line width=0.5pt]
  \def\u{11.88}
  % ---------------- left panel ----------------
  \begin{scope}
    \draw[->] (0,0) -- (0,96.9);
    \draw[->] (0,0) -- (85.0,0);
    \node[inner sep=0pt] at (-13.1,96.1) {\footnotesize$z$};
    \node[inner sep=0pt] at (93.8,-1.4) {\footnotesize$x$};
    \foreach \i in {0,...,6} {\fill (0,{\i*\u}) circle[radius=1.9pt];}
    \foreach \i in {0,...,5} {\fill ({\i*\u},0) circle[radius=1.9pt];}
    \draw (0,{6*\u}) -- ({5*\u},0);
  \end{scope}
  % ---------------- right panel ----------------
  \begin{scope}[xshift=166.3]
    \draw[->] (0,0) -- (0,96.9);
    \draw[->] (0,0) -- (85.0,0);
    \node[inner sep=0pt] at (-10.1,98.0) {\footnotesize$z$};
    \node[inner sep=0pt] at (95.3,0.1) {\footnotesize$x$};
    \foreach \i in {0,...,5} {\fill (0,{\i*\u}) circle[radius=1.9pt];}
    \foreach \i in {0,...,4} {\fill ({\i*\u},0) circle[radius=1.9pt];}
    \fill ({\u},{4*\u}) circle[radius=1.9pt];
    \draw (0,{5*\u}) -- ({\u},{4*\u}) -- ({4*\u},0);
  \end{scope}
\end{tikzpicture}

\caption{左：只有一个 block 的非正则奇点的牛顿多边形，原文图注在此把它说成情形 a 与 c。右：领头阶具有分割 $[N-1,1]$，原文图注在此把它说成情形 b。}
\label{fig:2}
\end{figure}

\typo{原文 p.11 的正文说「情形 a、b 用左图表示，情形 c 用右图表示」，而同一页图 2 的图注说的是「左图：只有一个 block 的非正则奇点，即情形 a 与 c；右图：领头阶具有分割 $[N-1,1]$，即情形 b」。两处对调了 b 与 c。对照图形可以判定：左图边界是一条直线（单个 block），右图边界折成两段（领头阶分成两块，即分割 $[N-1,1]$）。故正文正确、图注笔误。译文正文与图注均照录原文，此处以正文为准。}

非正则奇点的局部维数，也可以通过数「保持上述形式（即不改变特征值退化）的形变参数」求出。例如，若奇点是用一列杨图定义的（情形 d），那就只要数相应伴随轨道的维数，这个数文献 \cite{29} 已经给出。库仑分支维数等于模空间维数的一半。也可以用 Stokes 矩阵来数局部模空间的维数，这将在别处讨论 \cite{30}；这类 Stokes 矩阵对寻找相应场论的丛坐标很重要 \cite{30}。质量参数的个数从一阶极点矩阵的形式读出；耦合常数的个数则由对角形式中非正则奇点的参数个数得出。
